% GENERATED FILE. Edit canonical lessons, then run npm run generate:books.
% canonical-source-hash: fnv1a64:71a0ce8e6700e914
% canonical-lessons: MR-C167-svas-ghene, MR-C167-sodun-dene, MR-C167-adlabadal-karne, MR-C167-ghasarne, MR-C167-samsay-ghene

\chapter{To breathe, To give up, To exchange, To slip, To suspect}
\label{ch:mr-to-breathe-to-give-up-to-exchange-to-slip-to-suspect}
\hlchaptermodality{167}

\begin{chapteropening}
\textbf{By the end of this chapter:} \emph{I can say to breathe, to give up, to exchange, to slip and to suspect.}

Complete the last lesson of chapter 167: I can say to breathe, to give up, to exchange, to slip and to suspect.
\end{chapteropening}

\section[\texorpdfstring{\mr{श्वास} \mr{घेणे} (śvās gheṇe)}{श्वास घेणे (śvās gheṇe)}]{\mr{श्वास} \mr{घेणे} (śvās gheṇe) — to breathe}

\label{lesson:MR-C167-svas-ghene}



\begin{quote}
Before the new one: say the Marathi for to carry, then the Marathi for to hit; to kill.
\end{quote}

\subsection*{You'll want to know: \mr{श्वास} \mr{घेणे}}
\textbf{\mr{श्वास} \mr{घेणे}} — \emph{śvās gheṇe} — \textquotedblleft{}to breathe\textquotedblright{}.

\begin{grammarlens}[title={What you've built}]
One more word for this chapter.
\end{grammarlens}

\subsection*{Your turn}
\begin{itemize}
\raggedright
  \item \emph{Say it:} \emph{śvās gheṇe}
  \item \emph{Say it:} \emph{śvās gheṇe}, once more
  \item \emph{Say it:} say \emph{śvās gheṇe}
  \item \emph{Recall:} say the Marathi for to carry, then the Marathi for to hit; to kill, then say \emph{śvās gheṇe} again
\end{itemize}

\begin{tcolorbox}[breakable,colback=teal!4,colframe=teal!35!black,title={Before you move on}]
What does \mr{श्वास} \mr{घेणे} mean? (\textquotedblleft{}To breathe\textquotedblright{}.) Say it once more.
\end{tcolorbox}
\section[\texorpdfstring{\mr{सोडून} \mr{देणे} (soḍūn deṇe)}{सोडून देणे (soḍūn deṇe)}]{\mr{सोडून} \mr{देणे} (soḍūn deṇe) — to give up}

\label{lesson:MR-C167-sodun-dene}



\begin{quote}
Before the new one: say the Marathi for to hit; to kill, then the Marathi for to breathe.
\end{quote}

\subsection*{You'll want to know: \mr{सोडून} \mr{देणे}}
\textbf{\mr{सोडून} \mr{देणे}} — \emph{soḍūn deṇe} — \textquotedblleft{}to give up\textquotedblright{}.

\begin{grammarlens}[title={What you've built}]
One more word for this chapter.
\end{grammarlens}

\subsection*{Your turn}
\begin{itemize}
\raggedright
  \item \emph{Say it:} \emph{soḍūn deṇe}
  \item \emph{Say it:} \emph{soḍūn deṇe}, once more
  \item \emph{Say it:} say \emph{soḍūn deṇe}
  \item \emph{Recall:} say the Marathi for to hit; to kill, then the Marathi for to breathe, then say \emph{soḍūn deṇe} again
\end{itemize}

\begin{tcolorbox}[breakable,colback=teal!4,colframe=teal!35!black,title={Before you move on}]
What does \mr{सोडून} \mr{देणे} mean? (\textquotedblleft{}To give up\textquotedblright{}.) Say it once more.
\end{tcolorbox}
\section[\texorpdfstring{\mr{अदलाबदल} \mr{करणे} (adlābadal karṇe)}{अदलाबदल करणे (adlābadal karṇe)}]{\mr{अदलाबदल} \mr{करणे} (adlābadal karṇe) — to exchange}

\label{lesson:MR-C167-adlabadal-karne}



\begin{quote}
Before the new one: say the Marathi for to breathe, then the Marathi for to give up.
\end{quote}

\subsection*{You'll want to know: \mr{अदलाबदल} \mr{करणे}}
\textbf{\mr{अदलाबदल} \mr{करणे}} — \emph{adlābadal karṇe} — \textquotedblleft{}to exchange\textquotedblright{}.

\begin{grammarlens}[title={What you've built}]
One more word for this chapter.
\end{grammarlens}

\subsection*{Your turn}
\begin{itemize}
\raggedright
  \item \emph{Say it:} \emph{adlābadal karṇe}
  \item \emph{Say it:} \emph{adlābadal karṇe}, once more
  \item \emph{Say it:} say \emph{adlābadal karṇe}
  \item \emph{Recall:} say the Marathi for to breathe, then the Marathi for to give up, then say \emph{adlābadal karṇe} again
\end{itemize}

\begin{tcolorbox}[breakable,colback=teal!4,colframe=teal!35!black,title={Before you move on}]
What does \mr{अदलाबदल} \mr{करणे} mean? (\textquotedblleft{}To exchange\textquotedblright{}.) Say it once more.
\end{tcolorbox}
\section[\texorpdfstring{\mr{घसरणे} (ghasarṇe)}{घसरणे (ghasarṇe)}]{\mr{घसरणे} (ghasarṇe) — to slip}

\label{lesson:MR-C167-ghasarne}



\begin{quote}
Before the new one: say the Marathi for to give up, then the Marathi for to exchange.
\end{quote}

\subsection*{You'll want to know: \mr{घसरणे}}
\textbf{\mr{घसरणे}} — \emph{ghasarṇe} — \textquotedblleft{}to slip\textquotedblright{}.

\begin{grammarlens}[title={What you've built}]
One more word for this chapter.
\end{grammarlens}

\subsection*{Your turn}
\begin{itemize}
\raggedright
  \item \emph{Say it:} \emph{ghasarṇe}
  \item \emph{Say it:} \emph{ghasarṇe}, once more
  \item \emph{Say it:} say \emph{ghasarṇe}
  \item \emph{Recall:} say the Marathi for to give up, then the Marathi for to exchange, then say \emph{ghasarṇe} again
\end{itemize}

\begin{tcolorbox}[breakable,colback=teal!4,colframe=teal!35!black,title={Before you move on}]
What does \mr{घसरणे} mean? (\textquotedblleft{}To slip\textquotedblright{}.) Say it once more.
\end{tcolorbox}
\section[\texorpdfstring{\mr{संशय} \mr{घेणे} (saṁśay gheṇe)}{संशय घेणे (saṁśay gheṇe)}]{\mr{संशय} \mr{घेणे} (saṁśay gheṇe) — to suspect}

\label{lesson:MR-C167-samsay-ghene}



\begin{quote}
Before the new one: say the Marathi for to exchange, then the Marathi for to slip.
\end{quote}

\subsection*{You'll want to know: \mr{संशय} \mr{घेणे}}
\textbf{\mr{संशय} \mr{घेणे}} — \emph{saṁśay gheṇe} — \textquotedblleft{}to suspect\textquotedblright{}.

\begin{grammarlens}[title={What you've built}]
One more word for this chapter.
\end{grammarlens}

\subsection*{Your turn}
\begin{itemize}
\raggedright
  \item \emph{Say it:} \emph{saṁśay gheṇe}
  \item \emph{Say it:} \emph{saṁśay gheṇe}, once more
  \item \emph{Say it:} say \emph{saṁśay gheṇe}
  \item \emph{Recall:} say the Marathi for to exchange, then the Marathi for to slip, then say \emph{saṁśay gheṇe} again
\end{itemize}

\begin{tcolorbox}[breakable,colback=teal!4,colframe=teal!35!black,title={Before you move on}]
What does \mr{संशय} \mr{घेणे} mean? (\textquotedblleft{}To suspect\textquotedblright{}.) Say it once more.
\end{tcolorbox}
